%%
%% precntnt.tex - LaTeX2e thesis class
%%
%% Copyright (C) 2010-2021 Mathew Topper <damm_horse@yahoo.co.uk>
%%
%%
%%   ABOUT
%%
%% This is frontmatter prior to the contents for a Latex2e template which
%% corresponds to the regulations regarding layout of a thesis submitted
%% within the University of Edinburgh. 
%%
%% INPUT THIS FILE USING THE /makefrontmatter{} COMMAND OR THE FORMATTING
%% WON'T WORK PROPERLY

%%%% DEDICATION

\dedication{%
\begin{normalsize}To my grandfather Andr\'e ``Pop Clown'' Hodeir, who taught me to fill my life with pursuits that feel like play.\end{normalsize}\\[0.2cm]%
}

%%%% ABSTRACT
\abstract{
Reinforcement learning (RL) offers a powerful framework for learning adaptive control policies from interaction, but its practical deployment remains constrained by what this thesis terms the \textit{simulation bottleneck}. That is, effective learning and generalisation in RL tends to rely on bespoke simulators and access to sufficiently diverse training scenarios. In this thesis, we investigate both the causes of this bottleneck and algorithmic strategies for mitigating it. First, we study instance overfitting in deep RL, showing that the representations learned by modern actor-critic algorithms can make agents overly specialised to their training environments and drives poor generalisation. In particular, we demonstrate that separating actor and critic representations promotes useful specialisation, reduces harmful instance-specific encoding, and improves transfer. To support our study, we introduce mutual-information–based metrics that quantify specialisation and instance overfitting in the agents' representations. We show that the generalisation error of agents at deployment admits an upper bound determined by these metrics. %Complementing this analysis, we develop high-throughput simulation benchmarks, \hbox{PixelBrax} and \hbox{MEAL}, that make it possible to study generalisation and adaptation at scale under controlled distributions of training and evaluation instances.

The second part of the thesis develops generative approaches for relaxing, and ultimately bypassing, the simulation bottleneck. We introduce Data-Regularised Environment Design, a method that augments limited training sets with synthetic but distributionally grounded environment instances generated by a learned model, thereby improving zero-shot transfer without incurring the distribution shift common in prior unsupervised environment design methods. Finally, we move beyond hand-crafted simulators altogether through world models learned from offline interaction data. We present a world modelling framework that maintains a persistent latent 3D representation of the environment and uses it to generate action-conditioned visual experience with improved spatial memory, geometric consistency, and long-horizon stability. Together, these contributions advance a broader research programme aimed at making RL less dependent on manually engineered simulators, and closer to a scalable, data-driven paradigm for learning generalist policies.
}


%%%% LAY SUMMARY
\summary{
Reinforcement learning is a way of training computers to make decisions through trial and
error. Rather than being told exactly what to do, an artificial agent performs actions, observes
their consequences, and gradually learns which behaviours lead to better outcomes. However,
learning through unrestricted interaction in the real world can be slow, expensive, or unsafe.
Reinforcement learning agents are therefore often trained in computer simulations.

This creates what this thesis calls the \emph{simulation bottleneck}. Before an agent can learn,
a suitable simulator must be built, and it must contain enough varied training situations to
prepare the agent for deployment. Constructing these simulators and scenarios can require
substantial human expertise and engineering. The effort avoided by not manually programming
the agent's behaviour can therefore reappear in the design of the world in which it learns.

This thesis addresses the simulation bottleneck at three levels. First, it studies why agents
trained on a limited set of simulated environments often perform poorly in new ones. The results
show that agents can retain details that identify individual training environments instead of
learning information that transfers more broadly.

Second, the thesis introduces Data-Regularised Environment Design, or DRED, which generates
additional training environments from a limited initial dataset. Unlike methods that generate
situations simply because they are novel or difficult, DRED keeps them grounded in the
distribution of environments that the agent is expected to encounter. This expands training
diversity while avoiding unrealistic or irrelevant scenarios, leading to better generalisation.

Finally, the thesis investigates whether the simulator itself can be learned from recorded
experience. It introduces PERSIST, a learned world model that maintains an internal
three-dimensional representation of its environment. This persistent representation allows it
to remember previously visited places and generate more geometrically consistent and stable
visual experiences over long interactions. Together, these contributions reduce the need for
manually designed training scenarios and simulators, moving reinforcement learning towards a
more scalable, data-driven approach.
}

%%%% ACKNOWLEDGEMENTS

\acknowledgements{
First and foremost, I would like to express my deepest gratitude to my PhD co-advisors, David Abel and Christopher Lucas. Dave, thank you for taking me under your wing during a challenging period of my PhD. Your support, both during the thesis and beyond it, has meant more to me than I can easily express. I feel immensely fortunate to count you as both a mentor and a friend. Chris, thank you for your steadfast support throughout this journey. I cannot recall a single conversation with you that did not leave me with a new insight, no matter how niche or far removed the topic I brought to you. Your breadth of knowledge across computation and cognition made every conversation incredibly rewarding.

I am also deeply grateful to my three wonderful ``unofficial advisors'', Prakash Panangaden, Pablo Samuel Castro, and Kaixin Wang, for generously hosting me at Mila and at Microsoft Research, and for remaining closely involved in the research projects that eventually became central chapters of this thesis. I was equally fortunate to benefit from the mentorship and support of Tim Pearce, Steven McDonagh, and Stefano Albrecht.

My thanks also go to Michael Dennis and Alessandro Suglia for agreeing to serve as my examiners. Michael's work on Unsupervised Environment Design had a profound influence on the direction of my research at the beginning of the PhD. It is therefore particularly meaningful that you are serving as an examiner at its conclusion.

Thank you to Trevor McInroe and Thomas Walker for being such exceptional research partners and friends. I feel incredibly lucky to have found people with whom I could tackle challenging problems while having so much fun along the way. More than anything, I am grateful that our work together gave me the chance to get to know you both so well. It was also a privilege to work with such brilliant collaborators as Tristan Tomilin, James Doran, Shangmin Guo, and Vaidotas Simkus.

Many people, particularly members of the Agents and MARBLE research groups and regular attendees of the RL Reading Group and Generalisation Luncheons, provided invaluable feedback, insights, and ideas that helped shape this thesis. In no particular order, thank you to Kale-ab Tessera, Leonard Hinkledey, Lukas Schäfer, Chentian Jiang, Adam Jelley, Matthew Jackson, Nishanth Anand, Johan Obando-Ceron, Minqi Jiang, Michael Beukman, Logan Bhamidipaty, Ben Sanati, Raul Steleac, Bonan Zhao, Octave Mariotti, Ignacio Carlucho, Michael Gutmann, Reginald McLean, Callum Tilbury, Joseph Kelliher, Mhairi Dunion, Sabrina McCallum, Xuehui Yu, Atish Dixit, Alain Andres, Elliot Fosong, Arrasy Rahman, Cillian Brewitt, Filippos Christianos, Elle Miller, Balint Gyevnar, Georgios Papoudakis, Massimiliano Tamborski, and Riccardo Sarpietro.

I am grateful to the IPAB Institute and the Robotics and Autonomous Systems CDT for giving me the opportunity and financial support to pursue a PhD at the University of Edinburgh.

I would also like to thank the friends and mentors who first inspired me to pursue a PhD: Paul Bruce, Giorgio Geraci, Fernando Acero, and Ernesto Tripodi. To my old friends Arthur Pasquier, Raúl Arribas, Zeyad Sakr, Emilie and Kieran Pham-Downie and Antoine Volard, thank you for always taking such a genuine interest in what I was up to.

To Mom and Dad; to Isabelle, Aurélie, Emmanuel, Samantha, Françoise and Monica; to Charles-Antoine, Guillaume, Paola and Stéphane; and to my dear nieces, nephews, and cousins, Gabrielle, Paloma, Louise, Romane, Juliette, Artur, Lluc, Jules, Charlie and Gustave: thank you for surrounding me with so much love, encouragement, and support over the years.

Lastly, to Masha, my partner, wife, and best friend. Thank you for always believing in me; for convincing me that I had it in me to embark on this crazy journey; for looking after me through so many deadlines; and for patiently reading through even more drafts. My love and gratitude know no bounds.
}

% Dave, Chris
% Masha
% Pablo Samuel Castro, Prakash Panangaden, Kaixin Wang,
% Stefano Albrecht, Tim Pearce, Steven McDonagh
% **Trevor McInroe**, **Tom Walker**, Shangmin Guo,  James Doran, Tristan Tomilin
% Kale-ab Tessera, Leo Hinkledey, Elle Miller, Chentian Jiang, Adam Jelley, Ben (?), Raul (?), Bonan, Joe Kelliher, Mhairi Dunion(?), Sabrina McCallum,, Matthew Jackson, Xuehui Yu, Atish (?), Vaidotas Simkus, Alain Andres, Elliot Fosong, Arrasy Rahman(?), Ruaridh Williams(?), Cillian Brewitt, Filippos Christianos, Mhairi Dunion, Elliot Fosong, Balint Gyevnar, Georgios Papoudakis, Arrasy Rahman, Lukas Schäfer, Massimiliano Tamborski, Riccardo Sarpietro
% Arthur Pasquier, Antoine Volard, Zeyad X, Ra{u}l, Ignacio Carlucho, 
% Paul Bruce, Giorgio Geraci, Fernando Acero, Ernesto Tripodi
% Michael Dennis
% Mom, Dad, Isabelle, Aurelie, Emmanuel and their wonderful partners Charles-Antoine, Guillaume and Paola. My dear nieces and nephews Gabrielle, Paloma, Louise, Romanne, Juliette, Artur, Lluc and Jules.



%%%% DECLARATION

%% Use a custom declaration

% \declaration{I did it.}

%% Use the standard regulation declaration. Enter your
%% name for the signature line.

\standarddeclaration{Samuel Garcin}
